\documentclass[10pt]{article}
\usepackage[margin=0.65in]{geometry}
\usepackage{fontspec}
\setmainfont{Noto Serif}
\usepackage{booktabs,tabularx}
\usepackage{microtype}
\microtypesetup{nopatch=footnote}
\usepackage[hidelinks]{hyperref}
\setlength{\parindent}{0pt}
\setlength{\parskip}{3pt}
\setlength{\emergencystretch}{2em}
\title{Systematic-Review Reporting and Audit-Materials Statement}
\author{Anonymous for peer review}
\date{25 August 2026}
\begin{document}\maketitle\small

\section*{Reporting position}
The article is a systematic economic synthesis, not a meta-analysis. PRISMA 2020 is used as a benchmark for transparent reporting of search, selection, limitations and review materials. The manuscript does not claim a complete vendor-level PRISMA identification flow because the archived project does not preserve sufficiently reliable pre-deduplication result counts for each proprietary database. Those counts are not reconstructed or imputed after the fact. The 510-record post-deduplication object is therefore described as a \emph{broad discovery audit}, not as demonstrably high-recall.

\section*{Population and selection architecture}
\begin{tabularx}{\linewidth}{p{0.31\linewidth}X}
\toprule
Element & Final reporting state \\
\midrule
Deduplicated broad discovery audit & 510 records accumulated across bibliographic, frontier and citation-chasing routes.\\
Historical first pass & 437 retained / 73 excluded at that archived stage.\\
Working academic candidate pool & 444 after canonicalisation and reclassification.\\
Source-checked candidates & 299: 186 independently verified plus 113 earlier source/seed checked.\\
First-stage substantive eligibility & Direction-neutral: source-defensible economic mechanism, estimate, institutional analysis or theory within predefined governance/mechanism scope.\\
Second-stage active synthesis & 148 academic works performing distinct analytical work; 151 source-checked non-active candidates remain in the archived rowwise screening register, with aggregate disposition reported in the current bundle.\\
Policy/context layer & 3 sources, excluded from the academic evidence count.\\
Forensic layer & 27 studies, 69 source-location rows, 51 load-bearing claim IDs.\\
Synthesis unit & 118 proposition-ledger claims before evidence-dependence adjustment.\\
Direct causal evidence & 5 of 148 academic works are coded causal/quasi-experimental; policy rankings are correspondingly bounded.\\
\bottomrule
\end{tabularx}

\section*{Tension and boundary-condition synthesis}
The central matrix distinguishes two record types. Type I records are genuine sign or ranking reversals for the same or closely comparable policy object. Type II records are simultaneous institutional trade-offs in which effects can coexist because actors, margins or welfare objects differ. Each row exposes the evidential status of its boundary: \emph{Direct}, \emph{Model-imposed}, \emph{Cross-study inference}, or an explicit mixed classification. This prevents the review from manufacturing contradictions merely in order to reconcile them.

\section*{Evidence quality and independence}
Evidence class is not used as a quality score. Proposition-level credibility is assessed using method-specific validity dimensions, while evidence dependence is coded separately for shared datasets, shocks, benchmarks, models, primitives and conceptual lineages. Confidence labels---mechanism established, contextually identified, independently corroborated and open proposition---are assigned only after both credibility and dependence checks. The four numbered synthesis propositions display their confidence label, principal boundary, provenance and major falsifier in the main manuscript and machine-readable workbook.

\section*{Supporting audit materials}
The submission bundle includes the 148-work academic plus three-source policy/context evidence-role crosswalk, credibility framework, codebook, Type I/Type II tension matrix, four synthesis-proposition register, high-centrality seed list, disposition audit and an archive-register inventory. The inventory gives exact canonical v13.10 filenames and SHA-256 hashes for the working corpus, proposition ledger, dependence register, full-text extraction register, verification tranche and policy/context register. Older byte-frozen rows are not recreated from prose or generated summaries when the original file is retained in the project file library.

\section*{Search limitation}
The archived core discovery endpoint is 19 August 2026; a public frontier refresh through 24 August 2026 did not require reversal of a load-bearing proposition. The project preserves the post-deduplication discovery denominator and search families but not reliable raw pre-deduplication counts by vendor. Consequently, the submission uses an auditable synthesis-disposition flow rather than a retroactively completed PRISMA identification diagram. This is a fixed limitation of the historical log, not an unfinished screening stage.

\section*{Reference}
Page, M.J., McKenzie, J.E., Bossuyt, P.M., Boutron, I., Hoffmann, T.C., Mulrow, C.D., et al. (2021). The PRISMA 2020 statement: an updated guideline for reporting systematic reviews. \textit{BMJ}, 372, n71. \url{https://doi.org/10.1136/bmj.n71}.
\end{document}
